\section{A minimum exponent of two}
\label{sec:minimum-two}

Specializing the complement quintic to a minimum exponent of two
reduces its endpoint signs to positive polynomial expansions. One remaining
pair is handled by a rational interval, completing this minimum uniformly.

\begin{theorem}
\label{thm:minimum-two}
For distinct primes $p,q,r$ and all $b,c\geq1$, the graph
$G_{p^2q^br^c}$ is not Laplacian integral. Consequently every
three-prime exponent vector with minimum entry at most two gives a
nonintegral graph.
\end{theorem}

\begin{proof}
Use the complement quotient and quintic from Section~\ref{sec:unit}.
For $a=2$, $h_{2,b,c}(0)<0$. Putting $s=b+c$ and $t=bc$ gives
\[
h_{2,b,c}(1)=(s-13)t^2-(2s+6)t+3s^3+s^2-3s-1.
\]
By symmetry order $b\leq c$. A unit exponent is covered by
Theorem~\ref{thm:unit}; if $b=2$ or $b=c$, use
Theorem~\ref{thm:repeated}. It remains to consider $3\leq b<c$.

When $b=3$ and $c=5+v$, $v\geq0$, the exact expansion is
\[
h_{2,3,5+v}(1)=12v^3+112v^2+268v+120>0.
\]
When $b\geq4$ and $c>b$, write $b=4+u$ and $c=5+u+v$, with
$u,v\geq0$. Then
\begin{align*}
h_{2,4+u,5+u+v}(1)
={}&(u^2+8u+19)v^3\\
&+(4u^3+38u^2+128u+170)v^2\\
&+(5u^4+62u^3+271u^2+502u+368)v\\
&+2u^5+32u^4+190u^3+504u^2+552u+160>0.
\end{align*}
All coefficients are positive. In either region the opposite signs at
$0$ and $1$ give a root $\mu\in(0,1)$ and, by~\eqref{eq:complement}
and the universal-vertex lift, a noninteger graph eigenvalue in $(V-1,V)$.

The only remaining pair is $(b,c)=(3,4)$. Here
\[
h_{2,3,4}(x)=x^5-53x^4+953x^3-6523x^2+13134x-7560,
\]
and $h_{2,3,4}(1)=-48$ while
$h_{2,3,4}(3/2)=13437/32>0$. A root lies in $(1,3/2)$.
Since $V=58$, its lifted graph eigenvalue lies in $(113/2,57)$,
an interval containing no integer. This exhausts all positive $b,c$.
Combining with Theorem~\ref{thm:unit} proves the minimum-entry assertion.
\end{proof}
